% !TeX program = xelatex
\documentclass{article}
\usepackage{kotex}
\usepackage[a4paper,left=3cm, right=3cm, top=2cm, bottom=2cm]{geometry}
\usepackage{amsmath}
\usepackage{graphicx}
\usepackage{caption}
\usepackage{setspace}
\usepackage{xcolor}
\usepackage{titlesec}
\usepackage{amssymb}
\usepackage{tcolorbox}
\usepackage{wrapfig}
\usepackage{amsthm}
\usepackage{multirow}
\usepackage{array}
\usepackage{tikz}
\usepackage{longtable}
\usepackage{pgfplots}
\pgfplotsset{compat=1.18}

\title{1.3 Matrices and Matrix Operations}
\date{}
\author{}
\setstretch{1.3}

\titleformat{name=\section, numberless}
  {\normalfont\large\bfseries\color{blue}}
  {}
  {0pt}
  {}
\geometry{a4paper, margin=1in}

% Red subheading style (matches the textbook's red section headers)
\newcommand{\redheading}[1]{{\vspace{2mm}\noindent\Large\bfseries\color{red!55!black} #1\par}\vspace{2mm}}

\newtheoremstyle{mystyle}
  {}{}{\itshape}{}{\bfseries}{.}{.5em}{}
\theoremstyle{mystyle}

\newcommand{\deftitle}[1]{
\tcbset{colback=teal!4, colframe=teal!65!black, title=#1, fonttitle=\bfseries, boxrule=0.5mm, arc=3mm}}

\begin{document}
\maketitle

{\itshape\color{blue!45!black}
Rectangular arrays of numbers arise in many contexts beyond augmented matrices. This section studies matrices as objects in their own right, defining addition, subtraction, and multiplication on them.\par}

\redheading{Matrix Notation and Terminology}

Rectangular arrays occur in many contexts. For example, the table below records the hours a student spent on three subjects during a week:

\begin{center}
\begin{tabular}{|l|c|c|c|c|c|c|c|}
\hline
 & Mon. & Tues. & Wed. & Thurs. & Fri. & Sat. & Sun.\\
\hline
Math & 2 & 3 & 2 & 4 & 1 & 4 & 2\\
\hline
History & 0 & 3 & 1 & 4 & 3 & 2 & 2\\
\hline
Language & 4 & 1 & 3 & 1 & 0 & 0 & 2\\
\hline
\end{tabular}
\end{center}

Suppressing the headings leaves a rectangular array called a \textbf{matrix}:
\[
\begin{bmatrix}2&3&2&4&1&4&2\\0&3&1&4&3&2&2\\4&1&3&1&0&0&2\end{bmatrix}
\]

\begin{tcolorbox}[colback=teal!4, colframe=teal!65!black, title=Definition 1, fonttitle=\bfseries, boxrule=0.5mm, arc=3mm]
A \textbf{matrix} is a rectangular array of numbers. The numbers in the array are called the \textbf{entries} of the matrix.
\end{tcolorbox}

\subsection*{EXAMPLE 1 | Examples of Matrices}
\[
\begin{bmatrix}1&2\\3&0\\-1&4\end{bmatrix},\quad
\begin{bmatrix}2&1&0&-3\end{bmatrix},\quad
\begin{bmatrix}e&\pi&-\sqrt2\\0&\tfrac12&1\\0&0&0\end{bmatrix},\quad
\begin{bmatrix}1\\3\end{bmatrix},\quad
\begin{bmatrix}4\end{bmatrix}
\]

The \textbf{size} of a matrix is given as (rows) $\times$ (columns): the matrices above have sizes $3\times2$, $1\times4$, $3\times3$, $2\times1$, and $1\times1$. A matrix with one row (like the second above) is a \textbf{row vector} (or \textbf{row matrix}); one with one column (the fourth above) is a \textbf{column vector} (or \textbf{column matrix}); the fifth is both. We use capital letters for matrices and lowercase for numerical quantities (\textbf{scalars}, always real unless noted otherwise), e.g.\ $A=\begin{bmatrix}2&1&7\\3&4&2\end{bmatrix}$.

The entry in row $i$, column $j$ of $A$ is denoted $a_{ij}$, so a general $m\times n$ matrix is
\[
A=\begin{bmatrix}a_{11}&a_{12}&\cdots&a_{1n}\\a_{21}&a_{22}&\cdots&a_{2n}\\\vdots&\vdots& &\vdots\\a_{m1}&a_{m2}&\cdots&a_{mn}\end{bmatrix} \tag{1}
\]
written compactly as $A=[a_{ij}]_{m\times n}$ or $A=[a_{ij}]$. The entry in row $i$, column $j$ is also written $(A)_{ij}$, so $(A)_{ij}=a_{ij}$; e.g.\ for $A=\begin{bmatrix}2&-3\\7&0\end{bmatrix}$, $(A)_{11}=2,(A)_{12}=-3,(A)_{21}=7,(A)_{22}=0$.

Row and column vectors are usually denoted by boldface lowercase letters, needing only a single subscript:
\[
\mathbf{a}=[a_1\ \ a_2\ \ \cdots\ \ a_n]\qquad\text{and}\qquad \mathbf{b}=\begin{bmatrix}b_1\\b_2\\\vdots\\b_m\end{bmatrix}
\]

A matrix $A$ with $n$ rows and $n$ columns is a \textbf{square matrix of order $n$}; the entries $a_{11},a_{22},\dots,a_{nn}$ (boxed below) form its \textbf{main diagonal}:
\[
\begin{bmatrix}\boxed{a_{11}}&a_{12}&\cdots&a_{1n}\\a_{21}&\boxed{a_{22}}&\cdots&a_{2n}\\\vdots&\vdots& \ddots&\vdots\\a_{n1}&a_{n2}&\cdots&\boxed{a_{nn}}\end{bmatrix} \tag{2}
\]

\redheading{Operations on Matrices}

So far matrices have only abbreviated linear systems. We now develop an ``arithmetic of matrices'' -- addition, subtraction, and multiplication.

\begin{tcolorbox}[colback=teal!4, colframe=teal!65!black, title=Definition 2, fonttitle=\bfseries, boxrule=0.5mm, arc=3mm]
Two matrices are \textbf{equal} if they have the same size and their corresponding entries are equal.
\end{tcolorbox}

\subsection*{EXAMPLE 2 | Equality of Matrices}
For $A=\begin{bmatrix}2&1\\3&x\end{bmatrix}$, $B=\begin{bmatrix}2&1\\3&5\end{bmatrix}$, $C=\begin{bmatrix}2&1&0\\3&4&0\end{bmatrix}$: $A=B$ only if $x=5$; $A=C$ is impossible for any $x$, since they differ in size.

\begin{tcolorbox}[colback=teal!4, colframe=teal!65!black, title=Definition 3, fonttitle=\bfseries, boxrule=0.5mm, arc=3mm]
If $A$ and $B$ have the same size, the \textbf{sum} $A+B$ adds corresponding entries, and the \textbf{difference} $A-B$ subtracts them. Matrices of different sizes cannot be added or subtracted.
\end{tcolorbox}
In symbols, for $A=[a_{ij}]$, $B=[b_{ij}]$ of the same size:
\[
(A+B)_{ij}=a_{ij}+b_{ij}\qquad\text{and}\qquad (A-B)_{ij}=a_{ij}-b_{ij}
\]

\subsection*{EXAMPLE 3 | Addition and Subtraction}
For $A=\begin{bmatrix}2&1&0&3\\-1&0&2&4\\4&-2&7&0\end{bmatrix}$, $B=\begin{bmatrix}-4&3&5&1\\2&2&0&-1\\3&2&-4&5\end{bmatrix}$, $C=\begin{bmatrix}1&1\\2&2\end{bmatrix}$:
\[
A+B=\begin{bmatrix}-2&4&5&4\\1&2&2&3\\7&0&3&5\end{bmatrix},\qquad
A-B=\begin{bmatrix}6&-2&-5&2\\-3&-2&2&5\\1&-4&11&-5\end{bmatrix}
\]
$A+C,B+C,A-C,B-C$ are all undefined (size mismatch).

\begin{tcolorbox}[colback=teal!4, colframe=teal!65!black, title=Definition 4, fonttitle=\bfseries, boxrule=0.5mm, arc=3mm]
For any matrix $A$ and scalar $c$, the \textbf{product} $cA$ multiplies every entry of $A$ by $c$; $cA$ is a \textbf{scalar multiple} of $A$.
\end{tcolorbox}
In symbols, $(cA)_{ij}=c(A)_{ij}=ca_{ij}$.

\subsection*{EXAMPLE 4 | Scalar Multiples}
For $A=\begin{bmatrix}2&3&4\\1&3&1\end{bmatrix}$, $B=\begin{bmatrix}0&2&7\\-1&3&-5\end{bmatrix}$, $C=\begin{bmatrix}9&-6&3\\3&0&12\end{bmatrix}$:
\[
2A=\begin{bmatrix}4&6&8\\2&6&2\end{bmatrix},\qquad
(-1)B=\begin{bmatrix}0&-2&-7\\1&-3&5\end{bmatrix},\qquad
\tfrac13C=\begin{bmatrix}3&-2&1\\1&0&4\end{bmatrix}
\]
$(-1)B$ is usually written $-B$.

Multiplying matrices entry by entry (as with addition) turns out not to be useful. Instead:

\begin{tcolorbox}[colback=teal!4, colframe=teal!65!black, title=Definition 5, fonttitle=\bfseries, boxrule=0.5mm, arc=3mm]
If $A$ is $m\times r$ and $B$ is $r\times n$, the \textbf{product} $AB$ is the $m\times n$ matrix whose $(i,j)$ entry is found by multiplying corresponding entries of row $i$ of $A$ and column $j$ of $B$, then adding the products.
\end{tcolorbox}

\subsection*{EXAMPLE 5 | Multiplying Matrices}
For $A=\begin{bmatrix}1&2&4\\2&6&0\end{bmatrix}$ ($2\times3$) and $B=\begin{bmatrix}4&1&4&3\\0&-1&3&1\\2&7&5&2\end{bmatrix}$ ($3\times4$), $AB$ is $2\times4$. E.g., row 2 of $A$ with column 3 of $B$: $(2)(4)+(6)(3)+(0)(5)=26$; row 1 of $A$ with column 4 of $B$: $(1)(3)+(2)(1)+(4)(2)=13$. Computing all entries:
\[
AB=\begin{bmatrix}12&27&30&13\\8&-4&26&12\end{bmatrix}
\]

Matrix multiplication requires the number of columns of the first factor to equal the number of rows of the second; otherwise the product is undefined. Writing the sizes side by side, the inner numbers must match, and the outer numbers give the product's size:
\[
\underset{m\times r}{A}\quad\underset{r\times n}{B}\ \to\ \underset{m\times n}{AB} \tag{3}
\]

\subsection*{EXAMPLE 6 | Determining Whether a Product Is Defined}
For $A$ ($3\times4$), $B$ ($4\times7$), $C$ ($7\times3$): $AB$ is defined ($3\times7$); $BC$ is defined ($4\times3$); $CA$ is defined ($7\times4$); $AC$, $CB$, $BA$ are all undefined.

In general, for $A=[a_{ij}]$ ($m\times r$) and $B=[b_{ij}]$ ($r\times n$), the row-column rule gives
\[
(AB)_{ij}=a_{i1}b_{1j}+a_{i2}b_{2j}+a_{i3}b_{3j}+\cdots+a_{ir}b_{rj} \tag{5}
\]

\redheading{Partitioned Matrices}

A matrix can be \textbf{partitioned} into smaller \textbf{submatrices} by inserting rules between selected rows/columns. For a general $3\times4$ matrix $A$, three useful partitions are into four submatrices, into its three row vectors $\mathbf{r}_1,\mathbf{r}_2,\mathbf{r}_3$, or into its four column vectors $\mathbf{c}_1,\mathbf{c}_2,\mathbf{c}_3,\mathbf{c}_4$.

\redheading{Matrix Multiplication by Columns and by Rows}

Partitioning lets us extract a single row or column of $AB$ without computing the whole product:
\[
AB=A[\mathbf{b}_1\ \ \mathbf{b}_2\ \ \cdots\ \ \mathbf{b}_n]=[A\mathbf{b}_1\ \ A\mathbf{b}_2\ \ \cdots\ \ A\mathbf{b}_n] \tag{6}
\quad(AB\text{ column by column})
\]
\[
AB=\begin{bmatrix}\mathbf{a}_1\\\mathbf{a}_2\\\vdots\\\mathbf{a}_m\end{bmatrix}B=\begin{bmatrix}\mathbf{a}_1B\\\mathbf{a}_2B\\\vdots\\\mathbf{a}_mB\end{bmatrix} \tag{7}
\quad(AB\text{ row by row})
\]
In words:
\[
j\text{th column of }AB=A[j\text{th column of }B] \tag{8}
\]
\[
i\text{th row of }AB=[i\text{th row of }A]B \tag{9}
\]

\subsection*{EXAMPLE 7 | Example 5 Revisited}
For $A,B$ as in Example 5, the second column of $AB$ is
\[
\begin{bmatrix}1&2&4\\2&6&0\end{bmatrix}\begin{bmatrix}1\\-1\\7\end{bmatrix}=\begin{bmatrix}27\\-4\end{bmatrix}
\]
and the first row of $AB$ is $\begin{bmatrix}1&2&4\end{bmatrix}\begin{bmatrix}4&1&4&3\\0&-1&3&1\\2&7&5&2\end{bmatrix}=\begin{bmatrix}12&27&30&13\end{bmatrix}$.

\redheading{Matrix Products as Linear Combinations}

\begin{tcolorbox}[colback=teal!4, colframe=teal!65!black, title=Definition 6, fonttitle=\bfseries, boxrule=0.5mm, arc=3mm]
If $A_1,\dots,A_r$ are matrices of the same size and $c_1,\dots,c_r$ are scalars, then $c_1A_1+c_2A_2+\cdots+c_rA_r$ is called a \textbf{linear combination} of $A_1,\dots,A_r$ with \textbf{coefficients} $c_1,\dots,c_r$.
\end{tcolorbox}

For an $m\times n$ matrix $A$ and $n\times1$ vector $\mathbf{x}=[x_1\ \cdots\ x_n]^T$,
\[
A\mathbf{x}=\begin{bmatrix}a_{11}x_1+\cdots+a_{1n}x_n\\ \vdots \\ a_{m1}x_1+\cdots+a_{mn}x_n\end{bmatrix}
=x_1\begin{bmatrix}a_{11}\\\vdots\\a_{m1}\end{bmatrix}+x_2\begin{bmatrix}a_{12}\\\vdots\\a_{m2}\end{bmatrix}+\cdots+x_n\begin{bmatrix}a_{1n}\\\vdots\\a_{mn}\end{bmatrix} \tag{10}
\]
which proves:

\begin{tcolorbox}[colback=red!4, colframe=red!60!black, title=Theorem 1.3.1, boxrule=0.5mm, arc=3mm]
If $A$ is $m\times n$ and $\mathbf{x}$ is $n\times1$, then $A\mathbf{x}$ is a linear combination of the column vectors of $A$ with coefficients from $\mathbf{x}$.
\end{tcolorbox}

\subsection*{EXAMPLE 8 | Matrix Products as Linear Combinations}
\[
\begin{bmatrix}-1&3&2\\1&2&-3\\2&1&-2\end{bmatrix}\begin{bmatrix}2\\-1\\3\end{bmatrix}=\begin{bmatrix}1\\-9\\-3\end{bmatrix}
\quad\Longleftrightarrow\quad
2\begin{bmatrix}-1\\1\\2\end{bmatrix}-1\begin{bmatrix}3\\2\\1\end{bmatrix}+3\begin{bmatrix}2\\-3\\-2\end{bmatrix}=\begin{bmatrix}1\\-9\\-3\end{bmatrix}
\]

\subsection*{EXAMPLE 9 | Columns of a Product $AB$ as Linear Combinations}
From Example 5, $AB=\begin{bmatrix}1&2&4\\2&6&0\end{bmatrix}\begin{bmatrix}4&1&4&3\\0&-1&3&1\\2&7&5&2\end{bmatrix}=\begin{bmatrix}12&27&30&13\\8&-4&26&12\end{bmatrix}$. By (6) and Theorem 1.3.1, the $j$th column of $AB$ is the linear combination of $A$'s columns with coefficients from column $j$ of $B$:
\[
\begin{bmatrix}12\\8\end{bmatrix}=4\begin{bmatrix}1\\2\end{bmatrix}+0\begin{bmatrix}2\\6\end{bmatrix}+2\begin{bmatrix}4\\0\end{bmatrix},\quad
\begin{bmatrix}27\\-4\end{bmatrix}=\begin{bmatrix}1\\2\end{bmatrix}-\begin{bmatrix}2\\6\end{bmatrix}+7\begin{bmatrix}4\\0\end{bmatrix}
\]
\[
\begin{bmatrix}30\\26\end{bmatrix}=4\begin{bmatrix}1\\2\end{bmatrix}+3\begin{bmatrix}2\\6\end{bmatrix}+5\begin{bmatrix}4\\0\end{bmatrix},\quad
\begin{bmatrix}13\\12\end{bmatrix}=3\begin{bmatrix}1\\2\end{bmatrix}+\begin{bmatrix}2\\6\end{bmatrix}+2\begin{bmatrix}4\\0\end{bmatrix}
\]

\redheading{Column-Row Expansion}

If $m\times r$ matrix $A$ is partitioned into columns $\mathbf{c}_1,\dots,\mathbf{c}_r$ and $r\times n$ matrix $B$ into rows $\mathbf{r}_1,\dots,\mathbf{r}_r$, then each term $\mathbf{c}_i\mathbf{r}_i$ is $m\times n$, and
\[
AB=\mathbf{c}_1\mathbf{r}_1+\mathbf{c}_2\mathbf{r}_2+\cdots+\mathbf{c}_r\mathbf{r}_r \tag{11}
\]
called the \textbf{column-row expansion} of $AB$.

\subsection*{EXAMPLE 10 | Column-Row Expansion}
For $AB=\begin{bmatrix}1&3\\2&-1\end{bmatrix}\begin{bmatrix}2&0&4\\-3&5&1\end{bmatrix}$: $\mathbf{c}_1=\begin{bmatrix}1\\2\end{bmatrix},\mathbf{c}_2=\begin{bmatrix}3\\-1\end{bmatrix}$; $\mathbf{r}_1=[2\ 0\ 4],\mathbf{r}_2=[-3\ 5\ 1]$. So
\[
AB=\begin{bmatrix}1\\2\end{bmatrix}[2\ \ 0\ \ 4]+\begin{bmatrix}3\\-1\end{bmatrix}[-3\ \ 5\ \ 1]
=\begin{bmatrix}2&0&4\\4&0&8\end{bmatrix}+\begin{bmatrix}-9&15&3\\3&-5&-1\end{bmatrix}
=\begin{bmatrix}-7&15&7\\7&-5&7\end{bmatrix}
\]

\redheading{Summarizing Matrix Multiplication}

Five equivalent ways to compute a matrix product:
\begin{enumerate}
    \item Entry by entry (Definition 5)
    \item Row-column method (Formula (5))
    \item Column by column (Formula (6))
    \item Row by row (Formula (7))
    \item Column-row expansion (Formula (11))
\end{enumerate}

\redheading{Matrix Form of a Linear System}

For $m$ equations in $n$ unknowns, equal matrices means the system is equivalent to a single matrix equation:
\[
\begin{bmatrix}a_{11}x_1+\cdots+a_{1n}x_n\\\vdots\\a_{m1}x_1+\cdots+a_{mn}x_n\end{bmatrix}=\begin{bmatrix}b_1\\\vdots\\b_m\end{bmatrix}
\quad\Longleftrightarrow\quad
\begin{bmatrix}a_{11}&\cdots&a_{1n}\\\vdots& &\vdots\\a_{m1}&\cdots&a_{mn}\end{bmatrix}\begin{bmatrix}x_1\\\vdots\\x_n\end{bmatrix}=\begin{bmatrix}b_1\\\vdots\\b_m\end{bmatrix}
\]
Naming these matrices $A,\mathbf{x},\mathbf{b}$, the system becomes the single equation
\[
A\mathbf{x}=\mathbf{b}
\]
$A$ is the \textbf{coefficient matrix}; adjoining $\mathbf{b}$ as a last column gives the augmented matrix $[A\mid\mathbf{b}]$.

\redheading{Transpose of a Matrix}

We close with two operations that have no analog in ordinary arithmetic.

\begin{tcolorbox}[colback=teal!4, colframe=teal!65!black, title=Definition 7, fonttitle=\bfseries, boxrule=0.5mm, arc=3mm]
For any $m\times n$ matrix $A$, the \textbf{transpose} $A^T$ is the $n\times m$ matrix obtained by interchanging the rows and columns of $A$: the first column of $A^T$ is the first row of $A$, and so on.
\end{tcolorbox}

\subsection*{EXAMPLE 11 | Some Transposes}
\[
A=\begin{bmatrix}a_{11}&a_{12}&a_{13}&a_{14}\\a_{21}&a_{22}&a_{23}&a_{24}\\a_{31}&a_{32}&a_{33}&a_{34}\end{bmatrix}\!,\
B=\begin{bmatrix}2&3\\1&4\\5&6\end{bmatrix}\!,\ C=[1\ 3\ 5],\ D=[4]
\]
\[
A^T=\begin{bmatrix}a_{11}&a_{21}&a_{31}\\a_{12}&a_{22}&a_{32}\\a_{13}&a_{23}&a_{33}\\a_{14}&a_{24}&a_{34}\end{bmatrix}\!,\
B^T=\begin{bmatrix}2&1&5\\3&4&6\end{bmatrix}\!,\ C^T=\begin{bmatrix}1\\3\\5\end{bmatrix}\!,\ D^T=[4]
\]

The entry in row $i$, column $j$ of $A^T$ equals the entry in row $j$, column $i$ of $A$:
\[
(A^T)_{ij}=(A)_{ji} \tag{14}
\]
For a square matrix, $A^T$ can also be obtained by ``reflecting'' $A$ about its main diagonal -- interchanging entries symmetric about that diagonal:
\[
A=\begin{bmatrix}1&-2&4\\3&7&0\\-5&8&6\end{bmatrix}\ \longrightarrow\ A^T=\begin{bmatrix}1&3&-5\\-2&7&8\\4&0&6\end{bmatrix} \tag{15}
\]

\redheading{Trace of a Matrix}

\begin{tcolorbox}[colback=teal!4, colframe=teal!65!black, title=Definition 8, fonttitle=\bfseries, boxrule=0.5mm, arc=3mm]
For a square matrix $A$, the \textbf{trace} of $A$, written $\operatorname{tr}(A)$, is the sum of the entries on its main diagonal. The trace is undefined if $A$ is not square.
\end{tcolorbox}

\subsection*{EXAMPLE 12 | Trace}
For $A=\begin{bmatrix}a_{11}&a_{12}&a_{13}\\a_{21}&a_{22}&a_{23}\\a_{31}&a_{32}&a_{33}\end{bmatrix}$ and $B=\begin{bmatrix}-1&2&7&0\\3&5&-8&4\\1&2&7&-3\\4&-2&1&0\end{bmatrix}$:
\[
\operatorname{tr}(A)=a_{11}+a_{22}+a_{33},\qquad \operatorname{tr}(B)=-1+5+7+0=11
\]

\section*{Exercise Set 1.3}
\begin{enumerate}
    \item[1.] Suppose $A,B,C,D,E$ have sizes $(4\times5),(4\times5),(5\times2),(4\times2),(5\times4)$. Determine whether each expression is defined, and if so, give the size of the result.\\
    (a) $BA$\quad (b) $AB^T$\quad (c) $AC+D$\quad (d) $E(AC)$\quad (e) $A-3E^T$\quad (f) $E(5B+A)$

    \item[3.] Using $A=\begin{bmatrix}3&0\\-1&2\\1&1\end{bmatrix}$, $B=\begin{bmatrix}4&-1\\0&2\end{bmatrix}$, $C=\begin{bmatrix}1&4&2\\3&1&5\end{bmatrix}$, $D=\begin{bmatrix}1&5&2\\-1&0&1\\3&2&4\end{bmatrix}$, $E=\begin{bmatrix}6&1&3\\-1&1&2\\4&1&3\end{bmatrix}$, compute if defined:\\
    (a) $D+E$\ (b) $D-E$\ (c) $5A$\ (d) $-7C$\ (e) $2B-C$\ (f) $4E-2D$\\
    (g) $-3(D+2E)$\ (h) $A-A$\ (i) $\operatorname{tr}(D)$\ (j) $\operatorname{tr}(D-3E)$\ (k) $4\operatorname{tr}(7B)$\ (l) $\operatorname{tr}(A)$

    \item[4.] Using the matrices from Exercise 3, compute if defined:\\
    (a) $2A^T+C$\ (b) $D^T-E^T$\ (c) $(D-E)^T$\ (d) $B^T+5C^T$\ (e) $\tfrac12C^T-\tfrac14A$\ (f) $B-B^T$\\
    (g) $2E^T-3D^T$\ (h) $(2E^T-3D^T)^T$\ (i) $(CD)E$\ (j) $C(BA)$\ (k) $\operatorname{tr}(DE^T)$\ (l) $\operatorname{tr}(BC)$

    \item[7.] Using $A=\begin{bmatrix}3&-2&7\\6&5&4\\0&4&9\end{bmatrix}$, $B=\begin{bmatrix}6&-2&4\\0&1&3\\7&7&5\end{bmatrix}$, use the row or column method to find:\\
    (a) the first row of $AB$\quad (b) the third row of $AB$\quad (c) the second column of $AB$\\
    (d) the first column of $BA$\quad (e) the third row of $AA$\quad (f) the third column of $AA$

    \item[9.] Using $A,B$ from Exercise 7:\\
    (a) Express each column of $AA$ as a linear combination of the columns of $A$.\\
    (b) Express each column of $BB$ as a linear combination of the columns of $B$.

    \item[11.] Find matrices $A,\mathbf{x},\mathbf{b}$ expressing each system as $A\mathbf{x}=\mathbf{b}$, and write out the equation.\\
    (a) $2x_1-3x_2+5x_3=7,\ 9x_1-x_2+x_3=-1,\ x_1+5x_2+4x_3=0$\\
    (b) $4x_1-3x_3+x_4=1,\ 5x_1+x_2-8x_4=3,\ 2x_1-5x_2+9x_3-x_4=0,\ 3x_2-x_3+7x_4=2$

    \item[13.] Express each matrix equation as a system of linear equations.\\
    (a) $\begin{bmatrix}5&6&-7\\-1&-2&3\\0&4&-1\end{bmatrix}\begin{bmatrix}x_1\\x_2\\x_3\end{bmatrix}=\begin{bmatrix}2\\0\\3\end{bmatrix}$\qquad
    (b) $\begin{bmatrix}1&1&1\\2&3&0\\5&-3&-6\end{bmatrix}\begin{bmatrix}x\\y\\z\end{bmatrix}=\begin{bmatrix}2\\2\\-9\end{bmatrix}$

    \item[15.] Find all values of $k$, if any, satisfying
    \[
    \begin{bmatrix}k&1&1\end{bmatrix}\begin{bmatrix}1&1&0\\1&0&2\\0&2&-3\end{bmatrix}\begin{bmatrix}k\\1\\1\end{bmatrix}=0
    \]

    \item[17.] Use the column-row expansion of $AB$ to express the product as a sum of matrix products, for $A=\begin{bmatrix}4&-3\\2&-1\end{bmatrix}$, $B=\begin{bmatrix}0&1&2\\-2&3&1\end{bmatrix}$.

    \item[21.] For the linear system in Example 5 of Section 1.2, express the general solution obtained there as a linear combination of column vectors containing only numerical entries. [\emph{Suggestion:} write the general solution as one column vector, split it into column vectors each containing at most one parameter, then factor out the parameters.]

    \item[23.] Solve for $a,b,c,d$:
    \[
    \begin{bmatrix}a&3\\-1&a+b\end{bmatrix}=\begin{bmatrix}4&d-2c\\d+2c&-2\end{bmatrix}
    \]

    \item[25.] (a) Show that if $A$ has a row of zeros and $B$ is any matrix for which $AB$ is defined, then $AB$ also has a row of zeros.\\
    (b) State and justify a similar result for a column of zeros.

    \item[27.] How many $3\times3$ matrices $A$ satisfy $A\begin{bmatrix}x\\y\\z\end{bmatrix}=\begin{bmatrix}x+y\\x-y\\0\end{bmatrix}$ for \emph{all} $x,y,z$?

    \item[29.] A matrix $B$ is a \textbf{square root} of $A$ if $BB=A$.\\
    (a) Find two square roots of $A=\begin{bmatrix}2&2\\2&2\end{bmatrix}$.\\
    (b) How many different square roots can you find of $A=\begin{bmatrix}5&0\\0&9\end{bmatrix}$?\\
    (c) Do you think every $2\times2$ matrix has at least one square root? Explain.

    \item[30.] Let $0$ denote the $2\times2$ zero matrix.\\
    (a) Is there a $2\times2$ matrix $A\ne0$ with $AA=0$? Justify your answer.\\
    (b) Is there a $2\times2$ matrix $A\ne0$ with $AA=A$? Justify your answer.

    \item[32.] Find a $4\times4$ matrix $A=[a_{ij}]$ satisfying:\\
    (a) $a_{ij}=i+j$\qquad (b) $a_{ij}=i^{j-1}$\qquad
    (c) $a_{ij}=\begin{cases}1 & |i-j|>1\\ -1 & |i-j|\le1\end{cases}$

    \item[33.] Type I, II, III items cost \$1, \$2, \$3 each; the table gives the numbers purchased each month:\\
    \begin{center}
    \begin{tabular}{|l|c|c|c|}
    \hline
     & Type I & Type II & Type III\\
    \hline
    Jan. & 3 & 4 & 3\\
    \hline
    Feb. & 5 & 6 & 0\\
    \hline
    Mar. & 2 & 9 & 4\\
    \hline
    Apr. & 1 & 1 & 7\\
    \hline
    \end{tabular}
    \end{center}
    What does the following product represent?
    \[
    \begin{bmatrix}3&4&3\\5&6&0\\2&9&4\\1&1&7\end{bmatrix}\begin{bmatrix}1\\2\\3\end{bmatrix}
    \]

    \item[35.] \textbf{(Working with Proofs)} Prove: if $A,B$ are $n\times n$, then $\operatorname{tr}(A+B)=\operatorname{tr}(A)+\operatorname{tr}(B)$.
\end{enumerate}

\end{document}
